\section{Full three-dimensional normal cuts in affine cell families}
\label{sec:cut-cells}

The boundary constructor in Section~\ref{sec:cut-boundary} retains the complete
cut boundary and its chamber carriers. We now also retain the interior cells,
giving a full raw-cut CW description with oriented attachments in all four
dimensions. The representation uses binary affine families and includes the
original source charts and unscaled normal coordinates for reconstruction
of local three-ball prototypes. Describing the interior therefore does not
require expansion of a large normal vector.

The source remains a validated finite connected compact orientable
triangulation with one torus boundary. The normal surface may be empty,
disconnected, one-sided or a large multiple. The result describes
$M'=M\setminus\operatorname{int}\nu(F)$, with its boundary marked as a
subcomplex. It is a raw cut, without a parallelity collapse or an admissible
bounded handle core.

\subsection{All cell dimensions and their source charts}

The zero-cells are all original source vertices and the two cut vertices
at each normal intersection with a source edge. Normal vertices are boundary
vertices; original vertices retain their original boundary status. The
global edge orientation distinguishes the two cut sides, including loop edges.

The one-cells are both copies of every normal arc on a canonical source face
and all $w+1$ gap segments on every original edge of weight $w$. Interior
source edges are included. Gap endpoints split only at ranks zero and $w$;
an interior gap joins the upper copy of point $r-1$ to the lower copy of
point $r$. A normal-arc copy has affine maps to its cut-vertex endpoints.

The two-cells are both local copies of every normal disc and every region
of every canonical source face cut along normal arcs. Canonical faces include
one member of each glued pair and every original boundary face. Their centres
and corner arms retain the earlier cyclic recipes: triangles, quadrilaterals
and three-to-six-sided centres. Horizontal polygons and original boundary-face
regions are exactly the boundary two-cells.

The three-cells are all local chambers, with the established quad-chain and
triangle-arm numbering. Each is a three-ball in its source tetrahedron.
Chain ranks distinguish an empty-stack centre, the two quad endpoints and
interior quad prisms. Triangle-arm ranks distinguish the tip and triangle
prisms. The source row and rank determine the local prototype without
listing its copies.

Source charts are essential. Cellular boundary-matrix coefficients alone
need not specify a three-cell's topological attaching map. Every chamber
and incidence here also retains its source tetrahedron, local face or
normal-disc side and ordered rank. The packet contains the source face
permutations and normal vector, copied without aliasing mutable inputs.
These identify the abstract local boundary flags before source quotients.
The result is a reconstructible source cut, not merely a chain complex
with matching Euler characteristic.

\subsection{Oriented edge/polygon and polygon/chamber maps}

A normal-arc edge is incident to a source horizontal polygon, a target
horizontal polygon when the face is glued, and its adjacent original-face
region. The latter has inward rank $r$ or outward rank $r+1$, with the last
outward region replaced by the centre. Source and target triangle/quad
transitions retain separate incidence intervals.

A source-edge gap has an incidence for each occurrence on a canonical
source face. Its first corner arm, centre gap and last corner arm use at
most three affine pieces. Interior edge degree can be large. Keeping each
occurrence's own partition avoids a product of degree and common-refinement
size, preserving linear encoded storage.

Horizontal polygons attach once to their actual chamber with positive
outward incidence. Original boundary-face regions also attach once. An
interior original-face region has two incidences: positive on the canonical
owner and negative on its paired side. Their triangle/quad breakpoints
are retained independently. Two occurrences on the same chamber can cancel
as cellular coefficients; their source attaching data are still retained.

The source orientation determines polygon cyclic directions. Independent
replay uses the separate integer-midpoint polygon orientation model and
direct source-rank events. Singleton maps have a canonical slope-one
spelling; source side and permutation data remain separately recorded.

\begin{proposition}[Faithful raw-cut cells]
The four cell namespaces and attaching families describe exactly the CW
structure induced on $M'$ by the source normal cut. The marked cells are its
complete boundary subcomplex, and the signed cellular maps satisfy
\[
             \partial_1\partial_2=0,\qquad \partial_2\partial_3=0.
\]
All source self-identifications and local sides of one-sided surfaces remain.
\end{proposition}
\begin{proof}
Compatible triangle stacks and one quad stack cut each source tetrahedron
into the stated chain and arms of balls. Original faces cut into central
and corner polygons, original edges into ordered gaps, and every normal
disc, arc and edge point has two local frontier copies. These exhaust the
actual cut cells.

Normal matching and source permutations identify equal corner ranks and
their local edge and vertex flags across glued faces. Horizontal copies
meet the chamber on their stated local side. Each original-face region
meets its source and paired chambers at that same face rank. Thus all actual
attaching incidences, and only those incidences, are retained.

Each polygon traces its disc boundary, so oriented endpoints cancel. Each
local chamber has a sphere boundary, so its oriented polygon boundaries
cancel. These cancellations survive source identifications, even when two
occurrences identify on one chamber. Source charts retain the pre-quotient
flags and determine the full attaching maps, beyond their possibly cancelled
matrix entries. This proves both chain identities.

Horizontal frontier cells and their normal-arc and cut-vertex cells lie
on the new boundary. Original vertices, edge gaps and face regions lie
on the boundary exactly when they lie on the original boundary. Their marked
attachments reproduce Section~\ref{sec:cut-boundary}. Only local corner sides
were used, so one-sided monodromy is preserved without a global coorientation.
\end{proof}

\subsection{Family size and exact Euler accounting}

Let $V_M,E_M,F_M,T$ denote source cell counts, $D$ the normal disc count,
$A$ the normal arc count on canonical faces and $W$ the total global
edge weight. The cut has counts
\[
V'=V_M+2W,\quad E'=E_M+W+2A,\quad F'=F_M+A+2D,\quad C'=T+D.
\]
Hence $\chi(M')=\chi(M)+\chi(F)$, and its marked boundary has Euler
characteristic $2\chi(M')$. These exact binary conservation checks accompany
the source attaching-map proof; they do not replace it.

\begin{proposition}[Linear encoded raw-cut description]
Construction and independent source verification have polynomial bit
complexity in the source tetrahedron count $T$ and maximum coordinate bit
length $B$. The explicit integer/record accounting, including source charts,
uses $O(T(B+\log(T+1)))$ bits. No represented cell or chamber is enumerated.
\end{proposition}
\begin{proof}
There are at most $16T$ vertex, $30T$ edge, $26T$ polygon and $5T$ chamber
families. Each normal-arc family has at most two source horizontal, two
target horizontal and two original-region incidences, giving at most
$144T$ records. Source-edge occurrences add at most $36T$ records.
Gap endpoints use at most $18T$ pieces. Horizontal polygon/chamber maps
use at most $14T$ records; canonical and paired original-face maps add
at most $28T$ and $14T$, giving at most $56T$ in total. Cycles have bounded
length, and source reconstruction uses $O(T)$ face records and $7T$ coordinates.

All indices, counts and offsets have $O(B+\log(T+1))$ bits because
$D\le7T2^B$, $A\le4D$ and $W\le4D$. Each local occurrence has constantly
many rank events. Global sorting charges $O(T\log(T+1))$ record comparisons.
Source manifold and normal matching checks retain their polynomial cost.
Independent midpoint orientation uses fixed integer unit-tetrahedron
coordinates, with no floating approximation.

Separate incidence partitions avoid quadratic refinement at high-degree
edges. Exact arithmetic, source reconstruction and comparison are polynomial
in encoded size. The physical count $T+D$ is only a namespace endpoint,
never a native loop bound. The expanded finite chain oracle is a separate,
explicitly capped reference and is not an encoded-cost homology algorithm.
\end{proof}

\subsection{Source replay and the remaining hierarchy}

The APIs are \code{normal\_cut\_cell\_complex} and
\code{verify\_normal\_cut\_cell\_complex}. The schema is
\code{normal-cut-cells-v1}. It binds the fingerprint, reconstruction charts,
unscaled coordinates, every cell family, incidence sign, boundary flag,
cyclic recipe and chamber map. Invalid source geometry raises the existing
source error; malformed proofs return false. Cooperative cancellation
propagates without a partial result. No orbit discovery is started.

Literal finite checks build the three integral differentials and verify
both zero compositions, cyclic polygon closure, boundary-subcomplex closure,
opposite boundary-edge incidences and connected boundary vertex links.
Independent ranks over $\mathbb F_2$ and $\mathbb F_3$ are compared with
actual Regina cuts. The retained sphere-boundary nonball has Betti vectors
$(1,1,1,0)$ over $\mathbb F_2$ and $(1,0,0,0)$ over $\mathbb F_3$, detecting
its interior two-torsion. Homology alone still does not certify a ball.

The first supplied cut now has full raw interior data. Physical cell counts
can remain exponential in coordinate bits. General binary-cell homology,
parallelity removal, a bounded admissible core, arbitrary pattern transport,
general later-piece validation and the hierarchy depth bound are still
required. Ordinary recognition is unchanged. This supplies concrete compressed
cutting data for the original quasi-polynomial objective without claiming
that objective complete.

\subsection{Complete source and interior-homology evidence}

All 1,599 maintained tests pass with Regina available and no skips, and
all ten focused cut-cell tests pass. The complete audit covers 1,755 finite
cuts: 1,635 pilot surfaces, 48 mixed cases, 36 analytic examples and 36
relabellings. Every marked boundary genus inventory, Euler value and full
homology vector over $\mathbb F_2$ and $\mathbb F_3$ agrees with an actual
Regina cut. It checks 233,036 integral zero compositions, 638,563 edge/polygon
incidences and 260,842 polygon/chamber incidences. The boundary subcomplex
also passes its orientation and connected-vertex-link checks.

Ten large binary cases include meridians, vertex links and odd/even one-sided
multiples. In the 16,385-bit one-tetrahedron meridian case, seven vertex,
15 edge, 15 polygon and three chamber families retain the entire raw cut;
no represented cell is expanded. All 110 saved source proofs replay with
the cell constructor, producer orientation, summaries and orbit discovery
disabled. The audit's 682-hash source snapshot is retained at revision
\code{56a70f4da}. The sole subsequently changed snapshot file is the unused
benchmark driver; its archived Git content is authenticated separately while
all operation sources match the current implementation.

The actual full-cut homology oracle accounts for torsion in finite-field
homology using the universal coefficient theorem. If $H_1$ has $p$-torsion
rank $t_p$, then
\[
       b_1(\mathbb F_p)=\operatorname{rank}H_1+t_p,
       \qquad b_2(\mathbb F_p)=\operatorname{rank}H_2+t_p.
\]
Every cut component has nonempty boundary, so $b_3=0$. This oracle performs
explicit integer homology on expanded cuts; it is not the native constructor's
runtime. The retained audit took 1,743.062 seconds, with live sampling locating
its cost in Regina's homology calculation.

Reproduce the audit from \code{fast/} with:
\begin{quote}\small
\code{python -B -m normal\_orbit\_research.cut\_cells --output FILE}
\end{quote}
Its data and full-suite log are \code{data/cut-cells-audit.json} and
\code{data/cut-cells-tests.txt}. The native constructor and checker use
only the Python standard library; the independent oracle uses the project's
Regina extra, with Python distribution 7.4.1 and engine version string 7.4.

\subsection{Full raw-cut construction and serialization measurements}

Both measured arms construct a full representation of the same source cut.
Native constructs the binary CW packet, independently verifies its source
proof and serializes the entire result, including reconstruction charts and
coordinates. Regina constructs its explicit three-dimensional triangulation
and serializes its reversible labelled \code{tightEncoding}. Its exact
\code{tightDecoding} roundtrip is checked outside timing. Isomorphism
canonicalization is excluded. The two formats and geometric refinements
differ, while the native source proof authenticates its actual attaching maps.

Five shuffled paired rounds follow one warmup round, with an identical
Regina control. All 105 measured samples and 21 warmups finish. Euler values
agree in the timed cases; geometric validity comes from the separate source,
chain and homology audit, not that scalar timing check. Input preparation,
cleanup, garbage collection, decoding and comparison are outside timers.
The ratio is the median paired Regina/native ratio; values above one favour
native. The A/A column compares the identical Regina arms.
\input{tables/cut_cells.tex}

The 1,024- and 4,096-copy cases give ratios 5.407 and 25.792 with controls
0.962 and 0.997. The latter has native median 4.489 ms against 117.196 ms
Regina, with 41,869 serialized native bytes against 421,039 Regina bytes.
The primitive 16-tetrahedron case gives ratio 1.816 with control 1.045;
its native and Regina medians are 52.066 and 101.950 ms. Native is much
slower on tiny inputs: the one-copy ratio is 0.059. All regressions and
controls are retained. Neither the timing ratios nor the binary-family
bound imply a general recognition speedup.

Run the timing driver from \code{fast/}:
\begin{quote}\small
\code{python -B benchmark\_cut\_cells.py --output FILE}
\end{quote}
The current source-pinned records are \code{data/cut-cells-benchmark.json}.
The complete publication review replays the 110 saved proofs, verifies the
archived source snapshot and current operation/benchmark sources, checks
every timing round and reviews the article:
\begin{quote}\small
\code{python -B data/cut\_cells\_review.py --publication}
\end{quote}
Run that review from \code{synthesis/}. Parallelity removal, bounded core
geometry, pattern transport and the full hierarchy complexity proof remain
the next requirements for the original quasi-polynomial goal.
